\documentclass[12pt,a4paper]{article}
\usepackage{geometry}
\usepackage{tabto}
\usepackage{graphicx}
\usepackage{setspace}
\usepackage{longtable}
\usepackage{titlesec}
\usepackage{enumitem}
\usepackage{xcolor}

\geometry{margin=1in}
\setstretch{1.35}
\setlength{\parskip}{6pt}
\setlength{\parindent}{0pt}

\begin{document}


\titleformat{\section}
  {\large\bfseries}
  {\thesection}{1em}{}

\titleformat{\subsection}
  {\normalsize\bfseries}
  {\thesubsection}{1em}{}   


%================ TITLE PAGE ================
\begin{titlepage}
\centering

 \includegraphics[width=0.7\textwidth]{thapar_logo.jpg}
\rule{0pt}{3cm}

\vspace{1cm}

{\Large \textbf{THAPAR INSTITUTE OF ENGINEERING AND TECHNOLOGY}}\\
{\large Department of Computer Science and Engineering}

\vspace{1.5cm}

{\Huge \textbf{Project Proposal}}\\[0.5cm]
{\LARGE \textbf{Offline Hand Sign to Text Recognition System for American Sign Language}}

\vspace{1.5cm}

\textbf{Course Name \& Code:} UCS503 -- Software Engineering\\
\textbf{Degree \& Year:} B.Tech (3rd Year)\\
\textbf{Institute:} Thapar Institute of Engineering and Technology\\[0.5cm]

\textbf{Submitted By:}\\
Abhishek Srivastava \hfill - 1024030236\\
Keshav Aggarwal \hfill - 1024030234\\
Harshine Singh \hfill - 1024030232\\
Bhuvaneesh Singh Grewal \hfill - 1024030235\\[1cm]


\textbf{Academic Year:} 2026--27 ODD SEMESTER

\vfill

\end{titlepage}
%====================================================
%                    PROJECT PROPOSAL
%====================================================
\section{Higher-order goal}

This project aims to make basic sign-language communication accessible without depending on
paid cloud services, third-party APIs, or pretrained models. The broader motivation is
assistive and inclusive technology for the hearing- and speech-impaired community, but the
The immediate engineering objective is narrower and fully within the course scope: build an
end-to-end pipeline that recognizes static American Sign Language (ASL) hand signs, the 26
alphabet letters and 10 digits, from a live webcam feed, and converts them to on-screen text,
using only models trained from scratch on our own machines.


\section{Time-to-value}

\begin{itemize}[leftmargin=*]

  \item The pipeline is built in independently testable stages: raw-image pre-processing,
  model training, and live inference. Each stage produces a verifiable artifact before the
  next one begins.

  \item We prioritize fast feedback over perfection. The WebSocket link between the browser and
  backend was proven first with a trivial echo handler, before any machine-learning code was
  written, so integration bugs surface early rather than at the end.

  \item Success in the first iteration means a single captured frame correctly classified
  end-to-end. From there, we grow the class set and tighten accuracy.

\end{itemize}


\section{Problem Statement}

Hand-sign-to-text demos that students typically build tend to fall into two categories, and
Both are weak choices for a course project meant to show original engineering work.

API/cloud-dependent solutions call a hosted vision service for hand tracking or
classification. These need API keys, cost money or hit rate limits, send video data
off-device, and stop working offline.

Pretrained-model wrappers load an existing classifier and build a UI around it. That
shows integration skill, but not the ability to design, train, and evaluate a model from
scratch.

We want every learned component, the classifier that maps an image to a letter or digit, to
be trained by us, on our own data split, with no network calls at inference time. That means
figuring out how to structure a from-scratch CNN training pipeline, and how to split the work
across a three-person team so major components can be built in parallel, and how to serve a
locally trained model to a live browser feed at acceptable latency.


\section{Proposed Solution}

\subsection{Overview}

The system is a browser-to-server "Sign-to-Text" pipeline with three decoupled parts, each
owned mainly by one team member:

\begin{itemize}[leftmargin=*]

  \item Backend and CNN: a Java / Spring Boot service that exposes the frame-classification
  endpoint, holds the CNN training pipeline built with DeepLearning4j (DL4J), loads the
  trained model, and returns predictions.

  \item Frontend and Camera: a browser interface that accesses the webcam, captures frames,
  streams them to the backend, and shows the recognized letter or digit along with the
  accumulating text.

  \item Pre-processing: a Python pipeline that reads the raw ASL image dataset, cleans, and
  resizes it, and builds a reproducible train/test split for the CNN.

\end{itemize}

Testing runs across all three components rather than being confined to a separate module.


\subsection{Core workflow}

\begin{enumerate}[leftmargin=*]

  \item Pre-processing (offline, one-time): raw ASL images are read from the dataset's raw
  folder, not the provided "processed" folder, resized to $128 \times 128$, shuffled per
  class, and split 80/20 into our own train/test sets with a fixed random seed.

  \item Training (offline, one-time per iteration): a CNN is built and trained from scratch,
  no pretrained weights, on the processed training set, evaluated on the held-out test set
  After every epoch, keep the best checkpoint.

  \item Inference (live): the browser captures a frame and sends it to the backend, which runs
  It goes through the trained CNN and returns the predicted class.

  \item Display: the frontend shows the predicted letter or digit and appends it to a running
  text buffer, the user watches as it builds up in real time.

\end{enumerate}


\subsection{Operational constraints for an educational, from-scratch setting}

\begin{itemize}[leftmargin=*]

  \item No API keys, no third-party inference calls, no internet dependency at inference time.
  The trained model and server both run entirely on local hardware.

  \item No pretrained weights or transfer learning. The CNN's parameters are Xavier-initialized
  Moreover, learned entirely from our own training run.

  \item The dataset stays out of version control since it does not fit GitHub's practical size
  limits. Code lives on GitHub, the dataset separately, with its source and folder structure
  documented so every teammate can reproduce it locally.

\end{itemize}


\section{Solution Approach}

This is an engineering course project, so the focus is on a correct, reproducible,
from-scratch training-and-serving pipeline rather than inventing new model architectures.


\subsection{Dataset}

We use the ASL-HG (American Sign Language Hand Gesture) dataset: 36{,}000 images across 36
classes (A-Z and 0-9), collected from 10 volunteers, released with a raw-image folder and a
separate pre-processed folder. We pre-process the raw images ourselves, resizing, shuffling,
and splitting from scratch, instead of using the dataset's provided split, so every step from
raw pixels to the trained model is our own code.


\subsection{Model and training: from-scratch CNN in Java}

The classifier is a CNN with three convolution blocks (channels going $32 \to 64 \to 128$),
each followed by batch normalization, a non-linear activation, and max-pooling to halve the
spatial resolution, feeding into a small dense classification head over the 36 classes.
Instead of training in Python, we build and train the model directly in Java using
DeepLearning4j (DL4J), with ND4J handling the underlying tensor operations and DataVec
handling image loading and label inference from the folder structure. Training uses the Adam
optimizer and is evaluated after every epoch on the held-out test set, so the pipeline reports
a genuine train/test accuracy split instead of a single end-of-run number.


\subsection{Backend}

The Spring Boot backend exposes a lightweight endpoint that takes a captured frame from the
browser and returns the predicted label using the trained DL4J model loaded at startup.
Keeping the ML runtime and the serving layer in the same language avoids the need for a model export step.
(to ONNX, for example) that would otherwise be needed to run a Python-trained model outside
Python.


\subsection{Frontend and Camera}

The frontend is a browser interface using the browser's built-in \texttt{getUserMedia} camera
access to capture webcam frames. Captured frames go to the Spring Boot backend for
classification, and the predicted letter or digit is displayed to the user while recognized
characters accumulate into a running text buffer.

The frontend stays independent of the CNN implementation, talking to the backend through a
defined interface so the frontend and model components can be built and tested separately.


\subsection{CI and fast delivery enabler}

Each of the three components is independently buildable and testable, so a change to
Pre-processing does not require rebuilding the frontend, and changes to the CNN or backend don't
require touching the camera interface.

Python pre-processing checks, Java unit tests for model loading and prediction, and
Frontend integration tests catch pipeline breakage early, rather than at demonstration time.


\section{Evaluation Criterion}

Success is defined by measurable metrics tied to the core classification task.


\subsection{Primary evaluation metric}

Test-set classification accuracy: the percentage of held-out test images that the trained CNN
labels correctly.

\begin{itemize}[leftmargin=*]

  \item Target: $\geq 90\%$ top-1 accuracy on our own 20\% held-out split across all 36
  classes.

  \item Measured directly from the DL4J evaluation pass over the test iterator after training
  completes, logged per epoch to track convergence.

\end{itemize}


\subsection{Secondary evaluation metrics}

\begin{itemize}[leftmargin=*]

  \item Per-class precision/recall, to spot which visually similar signs (letters with
  near-identical hand shapes) The model confuses most.

  \item End-to-end inference latency, the time from frame capture in the browser to a
  displayed prediction, to confirm the system feels usable at a live, interactive pace.

  \item Reproducibility: re-running the pre-processing script with the fixed seed should
  produce an identical train/test split, checked by the test suite.

\end{itemize}


\subsection{Pilot validation plan}

Team members will each perform all 36 signs in front of the webcam under normal indoor conditions
lighting and record predicted-vs-actual labels, as a qualitative sanity check beyond the
held-out test accuracy.

We will also compare accuracy and confusion patterns between our from-scratch train/test split
and a quick baseline run on the dataset's own provided split, to confirm our pre-processing
It is not a quiet degradation of quality.


\section{Scalability}

\begin{enumerate}[leftmargin=*]

  \item Theoretical: the system is set up to grow from static alphabet/digit signs to dynamic,
  full-word gestures. Pre-processing, training, and serving stay decoupled, so a future
  A sequence-based classifier could replace the single-frame CNN without touching the frontend
  capture logic, and the backend endpoint is stateless per request, so it could later accept a
  A short buffered sequence of frames instead of one, for dynamic gesture recognition.

  \item Practical: the system runs entirely on commodity hardware. A laptop CPU is enough for
  Both training and inference at the current dataset scale, with no managed cloud service
  needed, which keeps it reproducible for every teammate and for grading.

\end{enumerate}


\section{Engine availability heuristic}

The tooling needed already exists and is freely available, so there is no need to build
infrastructure from scratch:

\begin{itemize}[leftmargin=*]

  \item Spring Boot for the backend web/REST layer and serving the trained CNN.

  \item DeepLearning4j, ND4J, and DataVec for CNN construction, training, and image data
  loading in Java.

  \item OpenCV and Pillow (Python) for raw-image reading and resizing during pre-processing.

  \item The browser's \texttt{MediaDevices}/\texttt{getUserMedia} APIs for webcam access, no
  An external SDK is required.

  \item Git/GitHub for source control, with the dataset hosted separately rather than
  committed to the repository.

\end{itemize}

Using these mature, freely available components lets us focus our effort on the from-scratch
model and pipeline design rather than rebuilding basic infrastructure that already exists.


\section{Project Scope and Deliverables}

\subsection{Initial Deliverable, First Iteration}

\begin{itemize}[leftmargin=*]

  \item Pre-processing: raw images resized, shuffled, and split 80/20 into train/test folders.

  \item CNN and Backend: CNN architecture and training loop in Java/DL4J, trained from scratch
  on the 36-class set, plus a Spring Boot endpoint that loads the trained model and classifies
  a single submitted frame.

  \item Frontend and Camera: a minimal browser frontend for camera access, frame capture, and
  display of the predicted letter/digit.

  \item Evaluation: per-epoch test-set evaluation, with the best checkpoint saved for serving.

\end{itemize}


\subsection{Subsequent Deliverables}

\begin{itemize}[leftmargin=*]

  \item A running text buffer in the frontend that accumulates recognized letters/digits into
  words.

  \item Confusion-matrix analysis and targeted data augmentation for the most-confused sign
  pairs.

  \item An automated test suite covering pre-processing determinism, model loading, backend
  prediction, and frontend-backend integration.

  \item Documentation explaining the Spring Boot wiring, DL4J training pipeline, pre-processing
  workflow, and frontend-camera integration.

  \item Stretch goal: extend from static single-frame classification to short buffered
  sequences, to support a small vocabulary of dynamic, full-word signs.

\end{itemize}


\section{Risks and Mitigations}

Training from scratch on 36 classes without pretrained weights may need more epochs or data
than a transfer-learning approach would. We are mitigating this by starting with a small,
well-understood architecture (3 conv blocks) and only adding capacity if we see underfitting.

The pipeline crosses languages, Python for pre-processing and Java for training, and a mismatch
in how images are resized, normalized, or labeled between the two could silently corrupt
training data. Tests check the output folder structure and image counts per class before training
is allowed to start.

Not every teammate starts with Java-web or DL4J experience. We are handling this with
plain-language documentation of new concepts as they come up, so owning a module doesn't
require re-deriving unfamiliar behavior in the framework from scratch.

The 36{,}000-image dataset is too large for practical Git version control, so it stays out of
the repository, with its public source and expected local folder layout documented in the
README.


\section{Summary}

This project is a fully self-trained, offline hand-sign-to-text system for the ASL alphabet
Furthermore, digits are built without pretrained models or third-party API dependencies. It hits the course
criteria through fast, independently verifiable pipeline stages, a clean three-way split of
work across Backend and CNN, Frontend and Camera, and Pre-processing, and a measurable
evaluation criterion centered on held-out classification accuracy.

It stays focused on disciplined, from-scratch engineering, our own data split, our own CNN,
our own training loop in Java, while leaving a clear path toward dynamic, full-word sign
recognition as future work.

\end{document}